% LỜI TỰA --- Quyển XIII
\chapter*{Lời Tựa}
\addcontentsline{toc}{chapter}{Lời Tựa}
\markboth{Lời Tựa}{Lời Tựa}

\begin{truyen}{Lời Tựa}{Questions Students Never Ask Aloud}
\chuhoa{G}{iáo viên bắt đầu bộ sách này với một câu hỏi đơn giản: làm thế nào để mỗi buổi học có thể để lại thứ gì đó đáng nhớ hơn là bài thi?}
\textit{(A teacher began this series with a simple question: how can each class leave behind something more memorable than a test?)}

Có những câu hỏi học trò giữ trong lòng mãi~--- vì không biết hỏi ai, hoặc sợ bị cho là ngớ ngẩn.
\textit{(Some questions students keep to themselves forever~--- because they don't know who to ask, or fear being seen as foolish.)}

Quyển mười ba chọn những câu hỏi đó ra và trả lời chúng một cách thẳng thắn, không vòng vo. Không phải mọi câu trả lời đều dễ nghe~--- nhưng tất cả đều thật.
\textit{(Volume thirteen selects those questions and answers them directly, without detour. Not every answer is easy to hear~--- but all of them are honest.)}

Bộ sách <<Truyện Tích Cảm Hứng Song Ngữ>> gồm 28 quyển, mỗi quyển một chủ đề riêng. Quyển XIII này mang chủ đề: \textbf{trả lời thẳng những câu hỏi bị giữ trong im lặng}.
\textit{(The series ``Inspiring Bilingual Tales'' contains 28 volumes, each with its own theme. Volume XIII focuses on: \textbf{trả lời thẳng những câu hỏi bị giữ trong im lặng}.})
\end{truyen}

\begin{baihoc}
Câu hỏi không được hỏi là câu hỏi không bao giờ được trả lời~--- và sự im lặng đó tốn nhiều hơn ta nghĩ.
\textit{(The question never asked is the question never answered~--- and that silence costs more than we think.)}
\end{baihoc}

\begin{ghinhoanh}
\textbf{Short English to remember:}

Brave questions deserve honest answers.

The unasked question is the most expensive silence.
\end{ghinhoanh}

\begin{tuhocnhanh}
1. Câu hỏi nào em đang giữ trong lòng mà chưa dám hỏi ai? \textit{(What question are you keeping inside that you haven't dared to ask?)}

2. Câu trả lời nào trong quyển này khiến em ngạc nhiên nhất? \textit{(Which answer in this volume surprised you most?)}

3. Điều gì ngăn cản em đặt câu hỏi trong lớp học? \textit{(What stops you from asking questions in class?)}
\end{tuhocnhanh}

\begin{center}
\textit{\color{lamporange}``Câu hỏi dũng cảm mở ra cánh cửa mà im lặng đã khóa lại.''}
\end{center}

\begin{flushright}
\textbf{Tôi Nguyễn Văn Sang}\\
\textit{GV THPT Nguyễn Hữu Cảnh - TP HCM}
\end{flushright}

\ngancach
